\section{基础能力与 Post-Training 的关系}

\subsection{两阶段范式}

Agentic Vision 揭示了一个重要的训练范式：

\begin{importantbox}{基础能力 + Post-Training 释放}
\begin{enumerate}
    \item \textbf{预训练阶段}：训练模型的\textbf{基础视觉能力}——检测、分割、识别。这些能力是"沉睡"的，模型已经知道如何做，但不知道何时做、如何组合。
    \item \textbf{Post-Training 阶段}：封装 Tool（如 Code Execution），通过强化学习激励模型形成\textbf{多轮工具调用}的行为模式，将基础能力"释放"出来。
\end{enumerate}
\end{importantbox}

这一范式对 AI Agent Research 有两层启发：

\begin{enumerate}
    \item \textbf{对 Prompting/Agent 设计者}：通过巧妙的 Visual Prompting 和 Tool 设计，可以在不改变模型参数的情况下，让之前不可行的任务变得可行
    \item \textbf{对模型训练者}：基础能力的质量决定了 Post-Training 后 Agent 的上限
\end{enumerate}

\subsection{Trigger Prompt 的角色}

在 Gemini 的官方 Demo 中，Prompt 中包含一些"触发词"来激发 Agentic Vision 能力：

\begin{itemize}
    \item "crop out all ..."
    \item "zoom in to see ..."
    \item "annotate on the image ..."
\end{itemize}

\begin{knowledgebox}{Trigger Prompt 是过渡性设计}
如果模型训练足够好，理论上不需要这些 Trigger Prompt——模型应该能够自主判断何时需要 zoom in、何时需要 annotate。Trigger Prompt 的存在说明当前模型的 Post-Training 还有优化空间。
\end{knowledgebox}

\subsection{本章小结}

Agentic Vision 的核心启示是：Agent 的能力上限由模型的基础能力决定，Post-Training 的作用是将这些基础能力组合、释放为完整的多轮推理行为。


